% !TEX root = ./Basilisk-celestialTwoBodyPoint-20190311.tex


\section{Test Description and Success Criteria}
The mathematics in this module are straight forward and can be tested in a series of input and output evaluation tests.

\subsection{Test 1}
The first test does a single-body celestial point. It places a spacecraft around Earth on an orbit with parameters given in \ref{tab:orbitParams}.
\begin{table}[htbp]
	\caption{Spacecraft Orbital Parameters}
	\label{tab:orbitParams}
	\centering \fontsize{10}{10}\selectfont
	\begin{tabular}{ c | c } % Column formatting,
		\hline\hline
		\textbf{Orbital Parameter}  & \textbf{Value}  \\
		\hline
		$a$    		        	     &  $2.8 R_{\text{earth}}	$   \\
		$e$           &  $0  $ \\
		$i$  & $0 \dg $\\
		$\Omega$  & $0\dg  $\\
		$\omega$  & $0\dg$\\
		$f$  & $60 \dg $\\
		\hline\hline
	\end{tabular}
\end{table}
The earth position and velocity vectors are both set to the zero vectors.

\subsection{Test 2}
The second test does the 2-body celestial point. It uses the same parameters at above and sets a second planet with zero velocity and position vector:

\begin{equation}
\leftexp{N}{\bm r_2} = \begin{bmatrix}500 & 500 &500 \end{bmatrix}^T (\text{m})
\end{equation}

\subsection{Finite-Difference Regression Tests}
Eight additional cases check the derivatives independently of the analytic truth model. The primary body starts at $[7,1,-2]^T\,10^6$ m with velocity $[-1,7,2]^T\,10^3$ m/s. The moving secondary body starts at $[-3,8,4]^T\,10^6$ m with velocity $[5,1,-3]^T\,10^3$ m/s. The spacecraft is at rest at the origin. Both bodies follow straight-line trajectories, consistent with the module's zero-relative-acceleration assumption.

The cases cover the moving secondary body, an absent secondary message, and secondary directions that remain parallel or antiparallel to the primary direction. Each configuration is evaluated at 0 s and 250 s, away from any fallback switching boundary. A central-difference step of 0.1 s is used.

The rate check differentiates the basis obtained from the output attitude and uses
\[
    \leftexp{N}{\bm\omega_{R/N}}
      = \frac{1}{2}\sum_{i=1}^{3}
        \leftexp{N}{\hat{\bm r}_i}\times\leftexp{N}{\dot{\hat{\bm r}}_i}.
\]
The acceleration check differentiates the published N-frame angular rate. Absolute tolerances are $10^{-11}$ rad/s for rate and $10^{-13}$ rad/s$^2$ for acceleration. These cases detect the component-frame mismatch from issue 1618 without copying the module's angular-acceleration equations.

\section{Test Parameters}
For the first two tests, the tested parameters are listed in Table~\ref{tab:errortol}.
\begin{table}[htbp]
	\caption{Error tolerance for each test.}
	\label{tab:errortol}
	\centering \fontsize{10}{10}\selectfont
	\begin{tabular}{ c | c } % Column formatting,
		\hline\hline
		\textbf{Output Value Tested}  & \textbf{Tolerated Error}  \\
		\hline
		$\bm\sigma_{R/N}$     		        	     & \input{AutoTeX/toleranceValue}  \\
		$\leftexp{N}{\bm\omega_{R/N}}$            &  \input{AutoTeX/toleranceValue} \\
		$\leftexp{N}{\dot{\bm \omega}_{R/N}}$   & \input{AutoTeX/toleranceValue}\\
		\hline\hline
	\end{tabular}
\end{table}




\section{Test Results}

The first two tests generate the following results. The finite-difference cases are checked by pytest:
\begin{table}[H]
	\caption{Test results}
	\label{tab:results}
	\centering \fontsize{10}{10}\selectfont
	\begin{tabular}{c | c  } % Column formatting,
		\hline\hline
		\textbf{Check} 						  		&\textbf{Pass/Fail} \\
		\hline
	   1.1	   			& \input{AutoTeX/passFail11} \\
	   1.2	   			& \input{AutoTeX/passFail12} \\
	   1.3	   			& \input{AutoTeX/passFail13} \\
	   2.1	   			& \input{AutoTeX/passFail21} \\
	   2.2	   			& \input{AutoTeX/passFail22} \\
	   2.3	   			& \input{AutoTeX/passFail23} \\
	   \hline\hline
	\end{tabular}
\end{table}
